We calculate the value of the equal payments as in the previous section:

\[R = \frac{A}{a\left(n,r\right)} = A\frac{r}{1-(1+r)^{-n}}.\]

For amortization we have the following new formulas which we will use to determine how much of each payment goes towards interest and how much goes towards the balance.

\[Interest = r\cdot Balance\]
\[Principal / repaid = payment - interest\]
\[Balance = old \ Balance - Principal \ repaid\]

We will use it to fill out an amortization schedule which is a table of all this information.

\begin{table}[]
    \centering
    \begin{tabular}{c|c|c|c|c}
         Pay \# & Pay & Interest & Principal Repaid & Balance\\
         \(0\) & x & x & x & \(A\) \\
         \(1\) & \(R\) & \(r\cdot A\) & \(R -rA\) & \(A - R -rA\) \\
         \vdots & \vdots & \vdots & \vdots & \vdots\\
         \(k - 1\) & \(R\) & ?  & ? & \(Ra\left(n-k+1,r\right)\) \\
         \(k\) & \(R\) & \(rRa\left(n-k+1,r\right)\) & \(R - rRa\left(n-k+1,r\right)\)  & \(Ra\left(n-k,r\right)\)
    \end{tabular}

\end{table}